\documentclass[10pt]{amsart}
\usepackage[a4paper,margin=22mm]{geometry}
\usepackage[no-math]{fontspec}
\setmainfont{Latin Modern Roman}[SmallCapsFont={lmromancaps10-regular.otf}]
\usepackage{amsmath,amssymb,amsthm,mathtools,bbm,hyperref,graphicx,etoolbox}
\emergencystretch=3em
\allowdisplaybreaks
\newtheorem{theo}{Theorem}[section]
\newtheorem{lemm}[theo]{Lemma}
\newtheorem{prop}[theo]{Proposition}
\newtheorem{coro}[theo]{Corollary}
\theoremstyle{definition}
\newtheorem{rema}[theo]{Remark}
\numberwithin{equation}{section}
\newcommand{\N}{\mathbb{N}}
\newcommand{\R}{\mathbb{R}}
\newcommand{\bbB}{\mathbb{B}}
\newcommand{\bbP}{\mathbb{P}}
\newcommand{\bbE}{\mathbb{E}}
\newcommand{\D}{\mathbb{D}}
\renewcommand{\S}{\mathbb{S}}
\renewcommand{\P}{\mathbb{P}}
\renewcommand{\H}{\mathbb{H}}

\DeclareMathOperator{\E}{\mathbb{E}}

\newcommand{\h}{\mathrm{h}}
\newcommand{\dd}{\mathrm{d}}
\DeclareMathOperator{\Hull}{Hull}
\DeclareMathOperator{\nee}{ne}
\DeclareMathOperator{\Ree}{Re}

\newcommand{\calS}{\mathcal{S}}
\newcommand{\calC}{\mathcal{C}}
\newcommand{\calD}{\mathcal{D}}
\newcommand{\calE}{\mathcal{E}}
\newcommand{\calG}{\mathcal{G}}
\newcommand{\calN}{\mathcal{N}}

\newcommand{\bfs}{\mathbf{s}}
\newcommand{\bfe}{\mathbf{e}}

\newcommand{\htec}{\mathbf{H}_{\mathrm{tech}}}
\newcommand{\ind}[1]{\mathbbm{1}_{\left\{#1\right\}}}

\renewcommand{\bar}[1]{\overline{#1}}
\renewcommand{\tilde}[1]{\widetilde{#1}}
\renewcommand{\epsilon}{\varepsilon}




%%%%%%%%%%%%%%%%%%%%%%%%%%%%%%%%%%%%%%%%%%%%%%%%%%%%%%%%%%%%%%%%%%%%%%%%%%%%%%%%%%%%%%%%%%%%%%%%

\graphicspath{{./figures/}}

\newcommand*{\mk}{\mkern -1mu}
\newcommand*{\Mk}{\mkern -2mu}
\newcommand*{\mK}{\mkern 1mu}
\newcommand*{\MK}{\mkern 2mu}

\hypersetup{urlcolor=purple, linkcolor=blue, citecolor=red}

\newcommand*{\relabel}{\renewcommand{\labelenumi}{(\theenumi)}}
\newcommand*{\romanenumi}{\renewcommand*{\theenumi}{\roman{enumi}}\relabel}
\newcommand*{\Romanenumi}{\renewcommand*{\theenumi}{\Roman{enumi}}\relabel}
\newcommand*{\alphenumi}{\renewcommand*{\theenumi}{\alph{enumi}}\relabel}
\newcommand*{\Alphenumi}{\renewcommand*{\theenumi}{\Alph{enumi}}\relabel}
\let\oldtilde\tilde
\renewcommand*{\tilde}[1]{\mathchoice{\widetilde{#1}}{\widetilde{#1}}{\oldtilde{#1}}{\oldtilde{#1}}}
\let\oldhat\hat
\renewcommand*{\hat}[1]{\mathchoice{\widehat{#1}}{\widehat{#1}}{\oldhat{#1}}{\oldhat{#1}}}
\let\oldforall\forall
\renewcommand*{\forall}{\mathrel{\oldforall}}
%%%%%%%%%%%%%%%%%%%%%%%%%%%%%%%%%%%%%%%%%%%%%%%%%%%%%%%%%%%%%%%%%%%%%%%%%%%%%%%%%%%%%%%%%%%%%%%%
\title[Sharp inverse-integral subordinator asymptotics]{Sharp stretched-exponential Laplace asymptotics for inverse-integral subordinators under regular variation and strong non-lattice hypotheses}
\author{B\'en\'edicte Haas}
\author{Bastien Mallein}
\date{Annales Henri Lebesgue8(2025),219–253; DOI10.5802/ahl.234}
\hypersetup{hidelinks}
\AtBeginDocument{\let\originalRef\ref\renewcommand{\ref}[1]{\ifstrequal{#1}{thm:asymp_distribution}{\href{https://doi.org/10.5802/ahl.234}{1.7 (source)}}{\ifstrequal{#1}{sec:proofscaling}{\href{https://doi.org/10.5802/ahl.234}{3.3 (source)}}{\originalRef{#1}}}}}
\title{Exact Laplace representation for a subordinator at inverse-integral time}
\author{Bénédicte Haas, Bastien Mallein}
\date{ahl:10.5802/ahl.234:published-2025; CC BY 4.0}
\begin{document}
\maketitle
\noindent\textit{Editorial note.} Only the designated primary theorem counts. Supporting named results and other statements are uncounted. Focused reading locates and understands the published author proof; it is not a stepwise independent proof verification.\par\medskip
\noindent\textit{Allowed standard prerequisites (editorial).}\begin{itemize}
\item Levy process fluctuation and strong Markov theory. Inverse exponent is a subordinator exponent and its Levy measure is the excursion lifetime measure. The PDE and resulting inverse-integral-time Laplace formula remain local obligations.
\item Fubini, dominated convergence and Laplace inversion. Basic exchange and inversion tools only; source proves relevant integrability, mixed derivatives and initial-value decomposition.
\end{itemize}\medskip
\subsection{Main results}\label{sec:def_fragmentations}

Motivated by the previous discussion, we consider now a generic subordinator $\xi$, with Laplace exponent
\[
\phi(q)=-\log \E\left[e^{-q \xi_1}\right] = \kappa + c q + \int_0^\infty (1 - e^{-qx}) \pi(\dd x)
\]
where $\kappa \geq 0$, $c \geq 0$ and $\pi$ is a measure on $(0,\infty)$ such that $\int_0^{\infty} (1 \wedge x) \pi(\dd x) < \infty$. We refer to $\pi$ as the Lévy measure of $\xi$, $c$ its drift and $\kappa$ its death rate. We define then the time change $\rho$ by
\[
\rho(t) := \inf \left\{ r \geq 0 : \int_0^r \xi_u \dd u \geq t \right\}, \quad t \geq 0,
\]
with the convention $\inf{\emptyset}=\infty$. Our first main result gives the following exact expression for the Laplace transform of $\xi_{\rho(t)}$, for any $t\geq 0$.



\begin{theo}\label{thm:lapTransform}
Let $\Phi(q)=\int_0^q \phi(s) \dd s$, $q \geq 0$ and $\Phi^{-1}: [0,\infty) \mapsto [0,\infty)$ be the inverse of $\Phi$. Then the function $\Phi^{-1}$ is the Laplace exponent of a subordinator, with a Lévy measure that we denote $L$, and for all $q > 0$ and $t \geq 0$, we have
\begin{equation}\label{eq:Laplace_intro}
\bbE\left[e^{-q\xi_{\rho(t)}}\right]~=~\phi(q)\int_0^{\infty} e^{-\Phi(q)x-\frac{t}{x}} xL(\dd x).
\end{equation}
\end{theo}

\section{The time-changed subordinator \texorpdfstring{$\xi_{\rho}$}{xiRho} and its Laplace exponent}\label{sec:Laplace}

In this and the next section, we work with $\xi$ a generic subordinator with Laplace exponent
\[
\phi(q)=\kappa+cq+\int_0^{\infty} \left(1-\exp(-qx)\right)\pi(\dd x), \quad q \geq 0
\]
where $\kappa \geq 0$, $c \geq 0$ and $\pi$ is a measure on $(0,\infty)$ such that $\int_0^{\infty} (1 \wedge x) \pi(\dd x)<\infty$. We assume throughout that $\pi(0,\infty)>0$ and $\xi_0=0$. We recall that $\rho(t)$ is the stopping time defined by
\begin{equation}\label{eq:time_change}
\rho(t)=\inf \left\{u \geq 0 : \int_0^u \xi_r \dd r >t \right\} \text{ for all }t\geq 0.
\end{equation}
Observe that $\rho(0) = \inf\{u \geq 0 : \xi_u > 0\}$ is either null almost surely, if $c > 0$ or $\pi((0,\infty)) = \infty$, or is strictly positive almost surely, when $c=0$ and $\pi$ is finite. In the latter case, it equals the first jump time of $\xi$, therefore $\xi_{\rho(0)}$ is equal to the first jump of $\xi$. Observe also that
\[
\int_0^{\rho(t)} \xi_r \dd r =t \quad \text{for all} \quad t <\int_0^{\bfe(\kappa)}\xi_r \dd r,
\]

\subsection{Definition and first properties of the measure \texorpdfstring{$L$}{L}}\label{sec:L}

We recall from the statement of Theorem~\ref{thm:lapTransform} that $\Phi$ is the primitive of $\phi$ satisfying $\Phi(0) = 0$, i.e.
\begin{align}\label{LK_Phi}
\Phi(q)&=\kappa q+\frac{c}{2}q^2+\int_{\R _+} \left(q-\frac{\left(1-\exp(-qx)\right)}{x}\right)\pi(\dd x) \\
\nonumber &= \kappa q+\frac{c}{2}q^2+ q \pi([1,\infty))+\int_0^{\infty} \left(qx \ind{x\,<\,1} -\left(1-\exp(-qx)\right)\right) \frac{\pi(\dd x)}{x}.
\end{align}
A key remark is that $\Phi$ is the Laplace transform of a Lévy process with diffusion coefficient $c$, drift $\kappa + \pi([1,\infty))$ and a Lévy measure given by the image of $\pi(\dd x)/x$ by $x \mapsto -x$. We denote $X$ such a process, satisfying
\begin{equation*}\label{def_X}
\bbE\left[e^{qX_t}\right]=e^{t \Phi(q)} \quad \text{for all } q,t \geq 0.
\end{equation*}
This class of processes, usually called \emph{spectrally negative Lévy processes}, has remarkable properties and has been studied extensively. We refer to Bertoin's book~\cite[Chapter~VII]{BLevy} for background. Here, we note that either $X$ oscillates ($\limsup_{t\rightarrow \infty} X_t=-\liminf_{t\rightarrow \infty} X_t=+\infty$) when $\kappa=0$, or tends to $+\infty$ ($\lim_{t\rightarrow \infty} X_t=+\infty$) when $\kappa>0$, using that $\Phi'(0)=\phi(0)=\kappa$. It is well known that $\Phi^{-1}:[0,\infty)\mapsto [0,\infty)$ is the Laplace exponent of a subordinator. More precisely,

\begin{lemm}[{\cite[Theorem~1, Chapter~VII]{BLevy}}]
The process $\sigma$ defined for $t \geq 0$ by
\[
\sigma_t = \inf\{ s > 0 : X_s > t\}=\inf\{ u >0:S_u>t\},
\]
where $S_t = \sup_{s\,\leq\,t} X_s$, is a subordinator with Laplace exponent $\Phi^{-1}$. Moreover $S$ is a local time at 0 of the reflected process $S-X$. Therefore, if we let $L$ denote the Lévy measure of $\sigma$, we have
\[
L(\dd x)=n(\zeta \in \dd x)
\]
where $n$ denotes the excursion measure away from 0 of $S-X$ and $\zeta$ the lifetime of an excursion.
\end{lemm}



Since $\Phi^{-1}(0)=0$ and, as $q\rightarrow \infty$, $\Phi^{-1}(q)/q$ converges to $(\kappa+\pi(0,\infty))^{-1}$ when $c=0$ and to $0$ otherwise, the subordinator $\sigma$ has no killing term and a drift equal to $\ind{c=0}(\kappa+\pi(0,\infty))^{-1}$. Consequently the function $\Phi^{-1}$ can be written for all $q\geq 0$ as
\begin{equation}\label{LK_invPhi}
\Phi^{-1}(q)=\frac{q \mathbbm 1_{\{c=0\}}}{\kappa+\pi(0,\infty)}+\int_0^{\infty} \left(1-\exp(-qx)\right)L(\dd x).
\end{equation}

\subsection{Proof of Theorem~\ref{thm:lapTransform}}\label{sec:proofThmain}

We now prove Theorem~\ref{thm:lapTransform}. In that aim we first introduce, for all $t \geq 0, q>0$,
\begin{equation}\label{def:f}
f(t,q):=(\Phi^{-1})'(q)\bbE\left[e^{-\Phi^{-1}(q) \xi_{\rho(t)}}\right]
\end{equation}
and observe that this bivariate function $f$ is the solution to a simple PDE.

\begin{lemm}\label{lem:easy}
The function $f$ is a strong solution to the equation
\begin{equation}\label{eq:easy}
\partial_q \partial_t f = \partial_t \partial_q f= f.
\end{equation}
Additionally, the function $t \in [0,\infty) \mapsto f(t,q)$ is $\calC^1$ on $(0,\infty)$ and continuous at 0, with $f(0,q)=\int_0^{\infty}xe^{-qx}L(\dd x)$, for all $q>0$.
\end{lemm}

In the proof of this lemma, we will note some regularity properties of $f$ that are a bit more precise than those stated here. Moreover, we will see further with Lemma~\ref{lem:laplaceLaplace} that $f$ is infinitely differentiable on $(0,\infty)^2$.

Instead of using this approach based on the infinitesimal generator, we prefer to use a more elementary method, which consists in rewriting the function $f$ as follows.
\begin{lemm}\label{lem:rewritten}
For all $t \geq 0$ and $q > 0$, the function $f$ defined in~\eqref{def:f} satisfies
\[
f(t,q) = \int_t^\infty \E\left[\frac{e^{-\Phi^{-1}(q)\xi_{\rho(s)}}}{\xi_{\rho(s)}} \right] \dd s.
\]
\end{lemm}


\begin{proof}
Fix $q>0, t \geq 0$.
We first observe that the change of variable $\rho(s) = u$ leads to
\[
\int_t^\infty \frac{e^{-\Phi^{-1}(q) \xi_{\rho(s)}}}{\xi_{\rho(s)}} \dd s = \int_{\rho(t)}^{\infty} e^{-\Phi^{-1}(q) \xi_u} \dd u.
\]
We underline that when $\kappa > 0$ the subordinator $\xi$ reaches $\infty$ in finite time and that the identity above is valid by using the convention $e^{-\infty} = 0$. Then by Fubini's theorem
\[
\int_t^\infty \E\left[\frac{e^{-\Phi^{-1}(q)\xi_{\rho(s)}}}{\xi_{\rho(s)}} \right] \dd s =
\E\left[\int_{\rho(t)}^{\infty} e^{-\Phi^{-1}(q) \xi_u} \dd u \right].
\]
Applying the strong Markov property for Lévy processes to the stopping time $\rho(t)$, we obtain
\begin{align*}
\E\left[\int_{\rho(t)}^{\infty} e^{-\Phi^{-1}(q) \xi_u} \dd u \right] &= \E\left[e^{-\Phi^{-1}(q) \xi_{\rho(t)}} \int_0^{\infty} e^{-\Phi^{-1}(q) \left(\xi_{u +\rho(t)} - \xi_{\rho(t)}\right)} \dd u \right]\\
&= \E\left[e^{-\Phi^{-1}(q) \xi_{\rho(t)}} \right] \int_0^\infty \E\left[e^{-\Phi^{-1}(q) \xi_{u}} \right] \dd u.
\end{align*}
Observing that $\E[e^{- \Phi^{-1}(q) \xi_t}] = e^{-t \phi(\Phi^{-1}(q))}=e^{-t \Phi'(\Phi^{-1}(q))}$, we conclude that
\[
\int_t^\infty \E\left[\frac{e^{-\Phi^{-1}(q)\xi_{\rho(s)}}}{\xi_{\rho(s)}} \right] \dd s = \frac{1}{\Phi' \circ \Phi^{-1}(q)}\E\left[e^{-\Phi^{-1}(q) \xi_{\rho(t)}} \right] = f(t,q). \qedhere
\]
\end{proof}


\begin{proof}[Proof of Lemma~\ref{lem:easy}]
Using that for all $a,b > 0$, $x \mapsto x^{a} e^{-bx}$ is bounded on $\R_+$, we obtain easily that the function $q \in (0,\infty) \mapsto f(t,q)$ is infinitely differentiable. Moreover, by Lemma~\ref{lem:rewritten}, for all $t \geq 0, q>0$,
\[
\int_t^\infty \E\left[\frac{e^{-\Phi^{-1}(q)\xi_{\rho(s)}}}{\xi_{\rho(s)}} \right] \dd s = f(t,q) < \infty,
\]
therefore $\E[e^{-\Phi^{-1}(q)\xi_{\rho(t)}}/\xi_{\rho(t)}] < \infty$ for almost all $t \geq 0$. Remarking that the function $t \mapsto e^{-\Phi^{-1}(q)\xi_{\rho(t)}}/\xi_{\rho(t)}$ is non-increasing, $t \mapsto \E[e^{-\Phi^{-1}(q)\xi_{\rho(t)}}/{\xi_{\rho(t)}}]$ is a non-increasing function. As a consequence, $t \in (0,\infty) \mapsto f(t,q)$ is a convex decreasing function and differentiable almost everywhere with
\[
\partial_t f(t,q) = -\E\left[\frac{e^{-\Phi^{-1}(q)\xi_{\rho(t)}}}{\xi_{\rho(t)}} \right] \quad \text{for almost all $t > 0$}.
\]
We then observe that for all $t,q > 0$, using Fubini's theorem
\begin{align*}
\int_q^\infty f(t,r) \dd r &= \int_q^\infty (\Phi^{-1})'(r)\E\left[e^{-\Phi^{-1}(r) \xi_{\rho(t)}}\right] \dd r \\
&= \E\left[\int_q^\infty (\Phi^{-1})'(r)e^{-\Phi^{-1}(r) \xi_{\rho(t)}} \dd r \right] = \E\left[\frac{e^{-\Phi^{-1}(q)\xi_{\rho(t)}}}{\xi_{\rho(t)}} \right].
\end{align*}
The function $t \in (0,\infty) \mapsto \E[{e^{-\Phi^{-1}(q)\xi_{\rho(t)}}}/{\xi_{\rho(t)}}]$ is therefore continuous, and consequently the function $t \in (0,\infty) \mapsto f(t,q)$ is $\calC^1$ for all $q > 0$ and
\[
\partial_t f(t,q) = -\int_q^\infty f(t,r) \dd r.
\]
Note that this implies that $q \in (0,\infty) \mapsto \partial_t f(t,q)$ is infinitely differentiable. Then, differentiating the above expression with respect to $q$ shows that $\partial_q \partial_t f(t,q) = f(t,q)$.

Similarly, by Fubini's theorem and the change of variables $\rho(s)=u$,
\begin{align*}
\int_t^\infty f(s,q) \dd s &= (\Phi^{-1})'(q)\bbE\left[\int_t^\infty e^{-\Phi^{-1}(q) \xi_{\rho(s)}} \dd s\right]\\
&= (\Phi^{-1})'(q) \E\left[\int_{\rho(t)}^{\infty} \xi_u e^{-\Phi^{-1}(q) \xi_u} \dd u \right]
\end{align*}
for $t,q>0$. Moreover, we note that for $u> 0$
\begin{multline*}
\E\left[\xi_{\rho(t)+u} e^{-\Phi^{-1}(q) \xi_{\rho(t)+u}}\right]\\
\begin{aligned}
&= \E\left[\left(\xi_{\rho(t)} +\left(\xi_{\rho(t) + u} - \xi_{\rho(t)}\right)\right) e^{-\Phi^{-1}(q) \xi_{\rho(t)}} e^{-\Phi^{-1}(q) \left(\xi_{\rho(t) + u} - \xi_{\rho(t)}\right)} \right]\\
&= \E\left[\xi_{\rho(t)} e^{-\Phi^{-1}(q)\xi_{\rho(t)}} \right] e^{-u \phi\left(\Phi^{-1}(q)\right)} + \E\left[e^{-\Phi^{-1}(q) \xi_{\rho(t)}}\right] u \phi'\left(\Phi^{-1}(q)\right) e^{-u \phi(\Phi^{-1}(q))}.
\end{aligned}
\end{multline*}
This leads to
\begin{multline*}
\int_t^\infty f(s,q) \dd s\\
\begin{aligned}
&= \frac{\left(\Phi^{-1}\right)'(q)}{\phi\left(\Phi^{-1}(q)\right)} \E\left[\xi_{\rho(t)} e^{-\Phi^{-1}(q)\xi_{\rho(t)}} \right] + \phi'\left(\Phi^{-1}(q)\right) \frac{\left(\Phi^{-1}\right)'(q)}{\phi\left(\Phi^{-1}(q)\right)^2} \E\left[e^{-\Phi^{-1}(q) \xi_{\rho(t)}} \right]\\
&= -\partial_q f(t,q),
\end{aligned}
\end{multline*}
since by definition of $f$
\[
\partial_q f(t,q) = (\Phi^{-1})''(q)\E\left[e^{-\Phi^{-1}(q) \xi_{\rho(t)}}\right] - (\Phi^{-1})'(q)^2 \E\left[\xi_{\rho(t)} e^{-\Phi^{-1}(q)\xi_{\rho(t)}} \right]
\]
and $(\Phi^{-1})' = 1/\phi (\Phi^{-1})$. Differentiating with respect to $t$ shows that we also have $\partial_t \partial_q f(t,q) = f(t,q)$.

The continuity at 0 of $t \mapsto f(t,q)$ is obvious, for all $q>0$, by using the expression of $f$ as an integral as noticed in Lemma~\ref{lem:rewritten}. It remains to identify the initial expression for $f(0,q)$ in terms of the measure $L$. Since $\xi_{\rho(0)}=0$ when $c>0$ or $\pi$ is infinite, and $\xi_{\rho(0)}$ is the first jump time of $\xi$ otherwise, one has
\[
\bbE\left[e^{-q \xi_{\rho(0)}}\right]=\left\{
\begin{array}{ll} 1 & \text{when $c>0$ or $\pi$ is infinite}\\ \frac{1}{\kappa+\pi(0,\infty)}\int_0^{\infty}e^{-qx} \pi(\dd x) & \text{when $c=0$ and $\pi$ is finite.}
\end{array}\right.
\]
When $c>0$ or $\pi$ is infinite, recalling~\eqref{LK_invPhi}, this indeed yields
\[
f(0,q)=(\Phi^{-1})'(q)=\int_0^{\infty} xe^{-qx} L(\dd x).
\]
When $c=0$ and $\pi$ is finite, we have
\[
f(0,q)=(\Phi^{-1})'(q) \frac{1}{\kappa+\pi(0,\infty)}\int_0^{\infty}e^{-\Phi^{-1}(q)x} \pi(\dd x).
\]
Using that in this case $\Phi(q)=(\kappa+\pi(0,\infty))q-\int_0^{\infty}(1-e^{-qx})\pi(\dd x)/x$, we see that
\[
q=(\kappa+\pi(0,\infty))\Phi^{-1}(q)-\int_0^{\infty}\left(1-e^{-\Phi^{-1}(q)x}\right)\pi(\dd x)/x,
\]
which, by differentiating, leads us to
\[
f(0,q)=(\Phi^{-1})'(q) \frac{1}{\kappa+\pi(0,\infty)}\int_0^{\infty}e^{-\Phi^{-1}(q)x} \pi(\dd x)=(\Phi^{-1})'(q) -\frac{1}{\kappa+\pi(0,\infty)}.
\]
By~\eqref{LK_invPhi} and since $c=0$, this last expression is equal to $\int_0^{\infty} xe^{-qx} L(\dd x)$ as expected.
\end{proof}

Solving the partial differential equation of Lemma~\ref{lem:easy}, we obtain an expression for the Laplace transform of the function $t \mapsto f(t,q)$ for each $q > 0$. This leads to an explicit expression of $f(t,q)$ in terms of the measure $L$, which proves Theorem~\ref{thm:lapTransform}.

\begin{lemm}\label{lem:laplaceLaplace}
For $\lambda>0$, we write $F_{\lambda} : q \in (0,\infty)\mapsto \int_0^{\infty} e^{-\lambda t} f(t,q) \dd t$. Then for all fixed $q>0$, we have
\begin{equation}\label{eq:Laplace}
F_{\lambda}(q)=\int_0^{\infty} \left(\int_0^{\infty} e^{-qx -\frac{t}{x}} xL(\dd x)\right) e^{-\lambda t} \dd t, \quad \forall \lambda>0
\end{equation}
which implies that $f(t,q)= \int_0^{\infty} e^{-qx -\frac{t}{x}} xL(\dd x)$ for all $t \geq 0, q>0$. In particular, $f$ is infinitely differentiable on $(0,\infty)^2$.
\end{lemm}


\begin{proof}
Fix $\lambda > 0$ and note that $F_{\lambda}(q)$ is finite for all $q > 0$, since the function $t \in (0,\infty) \mapsto f(t,q)$ is bounded. An integration by part gives
\begin{align*}
F_{\lambda}(q)=& \left[-\frac{e^{-\lambda t}}{\lambda}f(t,q) \right]_{t=0}^{t=\infty} + \frac{1}{\lambda} \int_0^{\infty}e^{-\lambda t} \partial _t f(t,q) \dd t \\
=& \frac{1}{\lambda} \int_0^{\infty} xe^{-qx} L(\dd x)+\frac{1}{\lambda} \int_0^{\infty}e^{-\lambda t} \partial _t f(t,q) \dd t,
\end{align*}
using that $e^{-\lambda t}f(t,q) \rightarrow 0 $ as $t\rightarrow \infty$, and that $f(t,q) \rightarrow \int_0^{\infty} xe^{-qx} L(\dd x)$ as $t \rightarrow 0$ by Lemma~\ref{lem:easy}. We then have, for all $q>0$,
\begin{align*}
\partial_q F_{\lambda}(q)&= \frac{-1}{\lambda} \int_0^{\infty} x^2e^{-qx} L(\dd x)+ \frac{1}{\lambda} \int_0^{\infty}e^{-\lambda t} \partial _q\partial _t f(t,q) \dd t \\
&= \frac{-1}{\lambda} \int_0^{\infty} x^2e^{-qx} L(\dd x)+ \frac{1}{\lambda} F_{\lambda}(q),
\end{align*}
using again Lemma~\ref{lem:easy} (we can apply Lebesgue's dominated convergence theorem since, for all $q_0>0$ and $t>0$,
\[
\sup_{q\,\geq\,q_0} \left|e^{-\lambda t}\partial_q \partial_t f(t,q)\right|=\sup_{q\,\geq\,q_0} \left|e^{-\lambda t} f(t,q)\right|=e^{-\lambda t} f(t,q_0),
\]
since $(\Phi^{-1})'$ is decreasing and $\Phi^{-1}$ is increasing).

The function $q \in (0,\infty) \mapsto F_{\lambda}(q)$ is therefore a solution to a linear differential equation of the first order. By standard arguments, it follows that
\[
F_{\lambda}(q)=e^{\frac{q}{\lambda}}\int_q^{\infty} e^{-\frac{u}{\lambda}} \left(\frac{1}{\lambda} \int_0^{\infty} x^2e^{-ux} L(\dd x) \right) \dd u.
\]
Using Fubini's theorem and the change of variable $t=(u-q)x/\lambda$, we then get the expression~\eqref{eq:Laplace} for all $q,\lambda>0$. The final equation follows from Laplace inversion theorem and the fact that the functions $t\mapsto f(t,q)$ and $t \mapsto \int_0^{\infty} e^{-qx -\frac{t}{x}} xL(\dd x)$ are continuous on $(0,\infty)$.
\end{proof}


\begin{proof}[{Proof of Theorem~\ref{thm:lapTransform}}]
Observe that the definition~\eqref{def:f} of $f$ implies that
\[
f(\Phi(q),t) = (\Phi^{-1})'(\Phi(q)) \E\left[e^{-q \xi_{\rho(t)}} \right] = \frac{1}{\phi(q)}\E\left[e^{-q \xi_{\rho(t)}} \right].
\]
As a result, Lemma~\ref{lem:laplaceLaplace} yields
\[
\bbE\left[e^{-q\xi_{\rho(t)}}\right] = \phi(q) f\left(t,\Phi(q)\right)=\phi(q) \int_0^{\infty} e^{-\Phi(q)x -\frac{t}{x}} \ xL(\dd x), \quad \forall t\geq 0, q>0,
\]
which completes the proof.
\end{proof}


\begin{thebibliography}{BGT87}
\bibitem[BGT87]{BGT} Bingham, Nick; Goldie, Charles; Teugels, Jozef Regular variation, Encyclopedia of Mathematics and Its Applications, 27, Cambridge University Press, 1987.
\bibitem[Ber98]{BLevy} Bertoin, Jean Lévy processes, Cambridge Tracts in Mathematics, 121, Cambridge University Press, 1998.
\bibitem[DR13]{DR13} Doney, Ron; Rivero, Victor Asymptotic behaviour of first passage time distributions for Lévy processes, Probab. Theory Relat. Fields, Volume 157 (2013) no. 1-2, pp. 1-45.
\bibitem[DR16]{DR16_err} Doney, Ron; Rivero, Victor Erratum to: “Asymptotic behaviour of first passage time distributions for Lévy processes”, Probab. Theory Relat. Fields, Volume 164 (2016) no. 3-4, pp. 1079-1083.
\bibitem[Kyp14]{KypFluctu} Andreas E. Kyprianou, Fluctuations of Lévy processes with applications. Introductory lectures, 2nd edition, Universitext, Springer,2014.
\end{thebibliography}
\end{document}
